\section{Scope, Evidence, and Working Assumptions}

\subsection{Purpose of this document}

The repository is not being treated as a toy order book bot project. The working target is a
serious discretionary-to-systematic research framework for intraday trading based on the same
information channels used by professional short-term prop traders: visible order book depth,
tape / prints, execution feedback, market context, and day-specific behavior. The immediate
goal of this document is therefore not to claim a finished strategy. The immediate goal is to
convert a diffuse knowledge base into a structured specification that can later be tested,
challenged, and encoded.

The materials mix beginner-oriented books, prop-infrastructure texts, and much more
practical lesson transcripts centered on reading the DOM. Because of that, the formalization
must separate three layers:

\begin{itemize}[leftmargin=1.4em]
  \item \term{Explicit claims}: statements that appear clearly and repeatedly in the materials.
  \item \term{Operational heuristics}: practical rules implied by repeated examples, even when not
        stated as formal laws.
  \item \term{Fuzzy discretionary content}: concepts that experienced traders clearly use, but which
        remain context-dependent and only partially specified.
\end{itemize}

\subsection{Sources covered}

This document is based on all relevant material present in the repository at the time of
writing.

\begin{longtable}{p{0.23\linewidth} p{0.51\linewidth} p{0.18\linewidth}}
\toprule
\textbf{Source group} & \textbf{What it contributes} & \textbf{Evidence weight} \\
\midrule
\endhead
Professional Scalping book & General scalping workflow, DOM and tape terminology, density bounce / breakout logic, cluster / PoC extensions, robot trading, and broad execution guidance. & Medium \\
Risk Management in Trading book & Deposit sizing, daily loss logic, per-trade limits, stop discipline, psychological traps, and journaling / review discipline. & Medium \\
Prop Trading on Moscow Exchange book & Prop infrastructure, risk-manager role, brokerage / exchange mechanics, stable access to market, server and operational controls. & Medium \\
Alexey Lipovoy lesson transcripts & Detailed discretionary setup logic from practice: densities, icebergs, robots, bounce, breakout, spring trades, plank trades, news reaction, market preparation, and self-analysis. & High \\
Existing repo data and viewer code & Concrete reminder that the project already contains old order-book-only Russian historical data and practical viewer tools, but still lacks a complete live capture stack for trades and executions. & High \\
\bottomrule
\end{longtable}

Two transcript files in the repository are empty in local storage at the time of writing:
``volatility harvesting'' and ``psychology.'' That does not mean those topics are absent from
the knowledge base overall; it means the local evidence for them is indirect and must be
reconstructed from the other materials.

\subsection{How to read confidence in the rest of the document}

Confidence in this document should be interpreted in a research sense, not in a trading
marketing sense.

\begin{itemize}[leftmargin=1.4em]
  \item \term{High confidence} means the idea appears repeatedly across the transcripts and is framed
        in a stable way.
  \item \term{Medium confidence} means the idea is clearly present but either broader than the raw
        evidence or partly inferred from examples.
  \item \term{Low confidence} means the idea is real enough to track, but still too fuzzy, too
        market-specific, or too weakly evidenced to treat as encoded truth.
\end{itemize}

\subsection{Project-level assumptions embedded here}

Several assumptions are not ``book knowledge'' in the narrow sense, but they are necessary to
make the repository coherent:

\begin{itemize}[leftmargin=1.4em]
  \item The long-term framework must use \term{order book + tape / prints + executions + context}, not
        order book snapshots alone.
  \item Historical Russian order-book-only data is useful for weak replay and directional sanity
        checks, but it is not enough to justify aggressive ML claims or fully trust backtest
        profitability.
  \item The correct progression is: formalize claims, test simple relationships, implement rule-based
        pattern detectors, then consider more adaptive filters and later ML overlays.
  \item Same-day calibration is a first-class research theme: the market should be observed in the
        opening phase to estimate what types of behavior are working \emph{today} for \emph{this}
        instrument before capital is deployed more aggressively.
\end{itemize}

\subsection{What this document is not}

This document is not a promise that every described setup still works unchanged in current
crypto or current Russian markets. It is not a set of guaranteed entry rules. It is not a
substitute for live observation. It is a structured bridge between discretionary training
materials and future empirical research.
